\scanpage{158}
\[
\begin{array}{lll}
Sl(2,\mathbb Z)&\rho,\sigma\ \mid&\rho^3=\sigma^2,\quad\sigma^4(=\rho^6)=1\\
&\varepsilon_0\ \varepsilon_1\ \mid&\deleted{(\varepsilon_1\varepsilon_0)^3\unsure{=}(\varepsilon_0\varepsilon_1\varepsilon_0)^2=1},\quad(\varepsilon_1\varepsilon_0)^6=1\\
&&\inserted{\varepsilon_0\varepsilon_1\varepsilon_0=\varepsilon_1\varepsilon_0\varepsilon_1}
\end{array}
\]
\[
\begin{array}{lll}
Gl(2,\mathbb Z)&\rho,\sigma,\unsure{\tau_0|r_0}\ \mid&\rho^3=\sigma^2,\ \sigma^4(=\rho^6)=1;\ r_0^2=1;\ r_0(\rho)=\rho^{-1};\ r_0(\sigma)=\sigma^{-1}\\
&r_0,r_\infty,r_{\unsure{1|a}}\ \mid&r_0^2=r_1^2=r_\infty^2=1\quad(r_1r_\infty)^2=(r_1r_0)^3,\quad(r_0r_\infty)^4=1\\
&r_0,\varepsilon_0,\varepsilon_1\ \mid&r_0^2=1,\quad r_0(\varepsilon_0)=\varepsilon_{\unsure{1|0}}^{-1},\quad r_0(\varepsilon_1)=\varepsilon_0^{-1},\\
&&\deleted{(\varepsilon_1\varepsilon_0)^3(\varepsilon_0\varepsilon_1\varepsilon_0)^{\unsure{2}}=1}\\
&&\inserted{\varepsilon_0\varepsilon_1\varepsilon_0=\varepsilon_1\varepsilon_0\varepsilon_1}\\
&&(\varepsilon_1\varepsilon_0)^6=1
\end{array}
\]

$\widetilde\varepsilon_0,\widetilde\varepsilon_1,\widetilde r_0$ \unsure{relèvent} du fait que $\varepsilon',\varepsilon'',r_0$ fixent le pt base $a_0$ de $\partial X$ ($\varepsilon',\varepsilon''$ sont $\equiv$ l'identité sur $\partial X$)

On définit ensuite.
\[
\widetilde\rho=(\widetilde\varepsilon_1\widetilde\varepsilon_0)^{-1}
=\widetilde\varepsilon_0^{-1}\widetilde\varepsilon_1^{-1}
\quad\text{i.e. }\widetilde\varepsilon_1\widetilde\varepsilon_0=\widetilde\rho^{-1}
\]
\margin{$\deleted{\unsure{\widetilde\rho(\widetilde\varepsilon_0)\Rightarrow}}$}
\[
\begin{aligned}
\widetilde\sigma&=\widetilde\varepsilon_0\widetilde\rho^{-1}&&\text{d'où}\quad\widetilde\varepsilon_0=\widetilde\sigma\widetilde\rho\\
&=\widetilde\varepsilon_0\widetilde\varepsilon_1\widetilde\varepsilon_0
&&\deleted{\unsure{\widetilde\varepsilon_1=\widetilde\rho\widetilde\sigma}}
\end{aligned}
\]
\deleted{\unsure{a-t-on}\quad$\widetilde\rho(\widetilde\varepsilon_0)=\widetilde\sigma(\widetilde\varepsilon_0)=$}

Or.
\[
\widetilde\sigma(\widetilde\varepsilon_0)=\sigma(\widetilde\varepsilon_0)=\widetilde\varepsilon_1
\]
d'où
\[
\widetilde\rho(\widetilde\varepsilon_0)=\widetilde\sigma^{-1}(\widetilde\varepsilon_0)=\sigma^{-1}(\widetilde\varepsilon_0)=\sigma(\widetilde\varepsilon_0)=\widetilde\varepsilon_1
\]
\margin{car $\widetilde\rho=\widetilde\sigma^{-1}\widetilde\varepsilon_0$\\
d'où $\widetilde\varepsilon_1=\widetilde\rho\widetilde\sigma$ et\\
$\widetilde\sigma=\widetilde\rho^{-1}\widetilde\varepsilon_1=\widetilde\varepsilon_1\widetilde\varepsilon_0\widetilde\varepsilon_1$}
\figure{\begin{tikzpicture}[x=1cm,y=1cm]
\draw (0,0) -- (6.4,0) .. controls (7.1,0) and (7.15,-.2) .. (7.15,-.9) .. controls (7.16,-1.6) and (7.35,-1.65) .. (7.75,-1.66) -- (10.4,-1.66);
\end{tikzpicture}}
\[
\widetilde\rho^{-3}=(\widetilde\varepsilon_1\widetilde\varepsilon_0)^3
\qquad\widetilde\sigma^2=(\widetilde\varepsilon_0\widetilde\varepsilon_1\widetilde\varepsilon_0)^2
\]
A-t-on $\underbrace{\widetilde\rho^{-3}=\widetilde\sigma^2}_{\widetilde\omega_0}$\quad i.e.
\[
\widetilde\varepsilon_1\widetilde\varepsilon_0\widetilde\varepsilon_{\unsure{1}}\widetilde\varepsilon_0\widetilde\varepsilon_1\widetilde\varepsilon_0
=\widetilde\varepsilon_0\widetilde\varepsilon_1\widetilde\varepsilon_0^2\widetilde\varepsilon_1\widetilde\varepsilon_0
\qquad ?
\]
\margin{oui}
On voudrait de plus que
\[
\widetilde\rho^{-6}=\widetilde\sigma^4
\]
soit le générateur canonique de $\operatorname{Ker}(S\mathcal T^+_{1,1}\longrightarrow\mathcal T_{1,1})$

\margin{158}
